% !PRINTED 667
%% [running head, not body text -- commented out 2026-08-07] Page furniture (running
%% head, folio, signature) is excluded from this edition by standing policy.
%% \S.\,184. \hfill \gesperrt{Case of an even $m$.} \hfill 667
\weberpage{667}
%<<CONT
If therefore we put, in accordance with (5),

\[
\mathfrak{E}^x \mathfrak{E}(r) = e(r)
\]

then it follows that

\weq{7}{\mathfrak{E}(r) = e(r)^m.}

But by (7) Theorem 2. is proved, and with it at the same time the
\gesperrt{whole Theorem I. for an odd $m$.}

This inference too fails for an even $m$, because the equation (6) is
then no longer solvable.

\begin{center}
\S.\,184.

\gesperrt{Preliminaries concerning the case of an even $m$.}
\end{center}

For the case in which $m$ is a power of 2 we take an entirely
different route, on which a few preliminary remarks may find their
place here first.

We have already seen in \S.\,182 that the first lemma is proved if we
can show that some power of $\varphi$ whose exponent is not divisible
by $q$ belongs to the principal class. In doing so, of the two
properties of the functional $\varphi$ only the first is used, namely
that the $m$th power of $\varphi$ is contained in the principal
class.

Now, by the general theorems in \S.\,154, we know in every field an
exponent $h$, namely the number of ideal classes, for which
$\varphi^h$ belongs to the principal class, whenever $\varphi$ is an
arbitrary functional in this field.

If, therefore, we can prove that the \gesperrt{class number $h$ in the field
$\Omega_m$, when $m$ is a power of 2, is an odd number,} then Lemma 1.
is thereby proved without anything further.

This is to be the aim of our considerations in the next chapters.

With regard to the second lemma the case is similar. Here one can
derive from the assumptions in \S.\,183, 2., quite as above, the
formula \S.\,183, (5); but from it one can conclude only that
$\mathfrak{E}(r)$ is the $\frac{1}{2} m$th power of a real unit
$e(r)$.

If, then, we assume the first lemma to be proved, the formula
\S.\,181, (14) yields:

\weq{1}{(r^n, x_0) = \varrho_n \sqrt{e(r^n)}\ \alpha_n\ \Pi(r^{m_2 n}, \eta),}
%CONT>>
